\section{From Effects to Their Ground}

The word ``metaphysics'' now needs rescuing from its caricature, because the caricature is what the instrumentalist refuted. In the caricature, the metaphysician closes his eyes to experience, consults a private faculty, and returns with news about a hidden world---assertion without discipline, precisely the wishful thinking the positivist hygiene was designed to kill. Whatever deserves that description deserves its fate. But the central tradition of Western metaphysics, and certainly the Catholic philosophical tradition, never worked that way, and said so with a bluntness that ought to embarrass the caricature.

Thomas Aquinas states the tradition's starting point in words an empiricist could sign: ``Our natural knowledge begins from sense. Hence our natural knowledge can go as far as it can be led by sensible things.''\endnote{Thomas Aquinas, \emph{Summa theologiae} I, q.~12, a.~12, corpus: ``naturalis nostra cognitio a sensu principium sumit, unde tantum se nostra naturalis cognitio extendere potest, inquantum manuduci potest per sensibilia.'' Latin checked 2026-07-24 at the Corpus Thomisticum online corpus; English from the public-domain translation of the Fathers of the English Dominican Province, checked the same day at New Advent. Work, question, and article govern over either website's presentation.} No innate blueprint of God, no private channel: the mind's material is what the senses deliver---the same stones, stars, and meter readings the sciences study. The question is only how far sensible things can \emph{lead}. And Aquinas's claim, argued rather than announced, is that they lead beyond themselves, because they present themselves as effects. His statement of the method is almost aggressively sober: ``from every effect the existence of its proper cause can be demonstrated, so long as its effects are better known to us\ldots\ Hence the existence of God, in so far as it is not self-evident to us, can be demonstrated from those of His effects which are known to us.''\endnote{\emph{Summa theologiae} I, q.~2, a.~2, corpus. Aquinas distinguishes demonstration \emph{propter quid} (through the cause, explaining why) from demonstration \emph{quia} (from the effect back to the cause's existence). Only the second is claimed for natural knowledge of God. Latin checked at Corpus Thomisticum, English at New Advent, 2026-07-24, as above.} Not seen in a vision; not felt in the heart---\emph{demonstrated}, the cold syllogistic word, from effects, the empiricist's own material.

What kind of reasoning is this? Not physics, and not in competition with physics. Physical explanation characteristically runs from state to state: the present configuration of a system is explained by an earlier configuration plus laws. Metaphysical inquiry asks a question that survives every such explanation, however complete: why is there a system in any configuration, and why do laws govern anything, rather than nothing at all? The contingent thing before me---this oak, this instrument, this galaxy---does not exist by necessity; it began, it depends, it could fail to be. Trace its dependence to prior things and the question does not shrink; it accumulates, because a chain of borrowers explains no loan. Either the regress of dependent things terminates in something that does not borrow its existence---something whose nature is not \emph{to be conditioned} but \emph{to be}---or the existence of the whole ensemble hangs on nothing. That, compressed past all its needed refinements, is the shape of the classical argumentation: not a laboratory report, but not an emotion either. It is inference from the manifest character of sensible things---their dependence, their contingency, their borrowed actuality---to what their existence requires. Whether it succeeds is a question of philosophical argument, and the arguments have been attacked and defended at the highest level of technical seriousness for centuries. What the arguments are \emph{not} is meaningless; that verdict required a criterion, and the criterion is dead.\endnote{This essay defends the admissibility and rational character of such reasoning; it deliberately does not execute any demonstration or claim that the reader is now in possession of one. The five ways of \emph{Summa theologiae} I, q.~2, a.~3, and their modern reconstructions and criticisms, are a literature of their own and lie outside this essay's scope. Nothing here argues from gaps in scientific explanation; the reasoning at issue begins from the existence and dependence of things, not from phenomena science has not yet explained---a distinction developed in the Scope appendix.}

Notice next how little is claimed, because the modesty is structural, not tactical. From effects, Aquinas insists, the mind cannot be led to see \emph{what God is}: ``the sensible effects of God do not equal the power of God as their cause,'' so that from them ``the whole power of God cannot be known; nor therefore can His essence be seen.'' What can be known is narrower and harder: ``we can be led from them so far as to know of God `whether He exists,' and to know of Him what must necessarily belong to Him, as the first cause of all things, exceeding all things caused by Him''---that all creatures stand to Him as caused to cause; that He is none of the things He causes; and that creatures fall short of Him not through any defect in Him but because He exceeds them all.\endnote{\emph{Summa theologiae} I, q.~12, a.~12, corpus, in the English Dominican translation as above. The doctrine of analogical rather than univocal naming of God, which disciplines all such language, is at I, q.~13, and is noted here without development.} This is a knowledge made mostly of negations and relations---a silhouette, not a portrait. The tradition that is accused of overreaching wrote the limits of its own reach into its founding texts.

It is worth pausing on how exactly this answers the instrumentalist's deepest instinct. His fear, remember, was assertion without discipline. But here is a discipline: premises drawn from public, sensible reality; inference by explicit steps, each attackable; conclusions bounded with almost legal care, claiming existence and relation while renouncing inner description. One may dispute a premise---many have---and the dispute will be philosophy, conducted with reasons. What one may not do, after the collapse of the criterion, is refuse the tribunal: to classify the argument as noise because its conclusion is not the sort of thing a detector could register. Of course it is not. The argument's entire point is that the detector, and every detectable thing, is an effect; a ground of all detectable things would be exactly what no detector could contain. To demand that the first cause show up as one more item inside the world's inventory is not rigor about God; it is a misunderstanding of what was ever being claimed.

The senses, then, are not despised in this tradition; they are trusted further than the empiricist trusts them---trusted to bear the weight of an inference that runs through them to what they cannot picture. The measurable world is not the rival of metaphysical reason but its raw material and its springboard. That conviction did not originate with Aquinas, and it is far older than the medieval schools. Before it was a philosophy it was a scriptural claim---a claim, remarkably, about what human reason can do at its best and what it is culpable for failing to do. The essay has argued on philosophy's ground to this point. It now reports, as the next witness, what the Catholic tradition has actually taught---because the teaching turns out to be the last thing an instrumentalist would predict a church to say.
